\section{Related Work}
\label{sec:related_work}

\subsection{Knowledge Graph Construction and Quality-Aware Extraction}

Traditional Knowledge Graph (KG) construction follows a systematic pipeline comprising information extraction, entity linking, relation mapping, and graph assembly \citep{ref_ji2021}. Early systems such as DBpedia \citep{ref_auer2007}, YAGO \citep{ref_suchanek2007}, and Freebase \citep{ref_bollacker2008} extract facts from structured and semi-structured sources. During the construction process, quality-aware mechanisms are actively employed to ensure data reliability. These mechanisms include confidence-based filtering, source voting, and entity disambiguation \citep{ref_dong2014}. Probabilistic knowledge fusion techniques integrate multiple extraction signals to assign reliability scores to generated triples, establishing a foundational layer of quality control directly within the extraction pipeline \citep{ref_dong2014}.

Embedding-based models map entities and relations into continuous vector spaces for link prediction and downstream inference \citep{ref_bordes2013,ref_wang2017}. Graph neural networks aggregate relational neighborhood information \citep{ref_schlichtkrull2018}. Surveys and position papers describe how LLMs can support KG construction, completion, and knowledge fusion \citep{ref_pan2023,ref_pan2025,ref_bian2025}. General prompting methods include chain-of-thought reasoning \citep{ref_wei2023} and reusable prompt patterns \citep{ref_white2023}; retrieval-augmented generation provides a mechanism for conditioning generation on external evidence \citep{ref_lewis2020}. These general methods should be distinguished from direct evaluations of Text-to-KG extraction. Public systems such as Neo4j Labs' LLM Knowledge Graph Builder expose a document-to-graph application with configurable schemas and an open-source extraction stack \citep{ref_neo4j_llm_graph_builder}. We include its reproducible core extractor as an external Text-to-KG baseline in Section~\ref{sec:eval:external_benchmark}.

\subsection{Knowledge Graph Quality Assessment}

Knowledge graph quality assessment characterizes graph reliability across multiple dimensions, including accuracy, completeness, consistency, and timeliness \citep{ref_zaveri2016, ref_paulheim2017}. Automated assessment methods evaluate these dimensions using diverse computational techniques. Statistical anomaly detection models value distributions to identify irregularities in numerical and literal data \citep{ref_wienand2014, ref_paulheim2014}. Embedding-based models evaluate triple plausibility by calculating distances in continuous vector spaces, providing a global semantic perspective on graph integrity \citep{ref_nickel2016}. Rule-based validation frameworks verify logical structures and schema compliance using formal constraint languages, ensuring explicit ontological consistency \citep{ref_kontokostas2014}. These methodologies provide comprehensive metrics for evaluating graph quality from local structural to global semantic perspectives.

Lin et al. introduce violation-inducing operations to construct controlled SHACL repair tasks and compare LLM-based repair strategies \citep{ref_lin2025}. Their results highlight the value of relevant constraints and local graph context in repair prompts. This is evidence about constraint-aware repair, rather than a validation of a general-purpose semantic judge. Crowdsourcing complements automated assessment for ambiguous quality judgments \citep{ref_acosta2013}. Broader surveys organize quality control around multiple dimensions and stages of KG construction and refinement \citep{ref_wang2021,ref_paulheim2017}. Our diagnostic report records executable structural and source-support checks within a single document graph.

\subsection{Knowledge Graph Quality Enhancement and Tool Orchestration}

KG enhancement techniques address missing facts and erroneous assertions. Completion methods based on relational paths \citep{ref_lao2011}, embeddings \citep{ref_bordes2013,ref_yang2015,ref_shi2017}, and neural architectures \citep{ref_dettmers2018} predict missing links. MINERVA learns a policy for navigating a KG to answer queries \citep{ref_das2018}, while MetaR addresses few-shot link prediction \citep{ref_chen2019}. These tasks provide mechanisms for inference, but are distinct from evaluating whether an edit preserves valid facts in a corrupted graph. AMIE mines positive logical rules \citep{ref_galarraga2013}, and RuDiK discovers positive and negative rules \citep{ref_ortona2018}. Other approaches incorporate logical rules or simple constraints into embedding objectives \citep{ref_guo2016,ref_ding2018}.

Neural and probabilistic approaches further combine logical structure with learned predictions. RNNLogic jointly trains a rule generator and a reasoning predictor through an EM-based procedure \citep{ref_qu2021}. NBFNet learns path-based representations through a generalized Bellman--Ford formulation for link prediction \citep{ref_zhu2021}; it is not an explicit rule-mining or KG-repair evaluation system. pLogicNet combines first-order rules in a Markov logic network with embedding-based variational inference \citep{ref_qu2019}. These methods motivate the use of structural priors alongside learned predictions. Our evaluated pipeline normalizes document-field triples and validates generated candidates. We separately audit a trial-graph optimizer and do not attribute pipeline performance to that optimizer.

LLM-based KG repair has been evaluated using explicit SHACL violations and alternative prompting strategies \citep{ref_lin2025}. More generally, ReAct demonstrates how language-model reasoning can be interleaved with actions to gather information from external sources \citep{ref_yao2023react}. This interaction pattern informs our ReAct-style baseline. Our study isolates diagnostic prompting and source-grounded candidate filtering using shared instructions and frozen candidate responses. A five-part SHACL-context baseline adapts Lin et al.'s design to the same document-field task and model budget. The empirical comparison, rather than the presence of prompting or constraints alone, determines which mechanisms are supported.
